\section{The Halting Horizon}

A repeatable trial must end, or it is not a trial. This does not mean
that every question posed to measurement receives a final answer. It
means that the record must know how much time has been spent and what
has been obtained by spending it. The grammar therefore introduces a
horizon. Observation is bounded by the time the apparatus can spend.

The horizon is often mistaken for weakness. If only the trial could run
longer, if only the resolution could be finer, if only the ledger could
include more intermediate cases, then perhaps the measurement would
become complete. The grammar replies that longer, finer, and more are
all further readings. They may be permitted, but they are not already
present. A bound is not the enemy of measurement. A bound is what lets a
reading be finished enough to be compared.

The halting horizon is the point at which finite approach becomes
finite observation. A residue may still remain. A future trial may shrink
it. A different gauge may read a feature that this gauge cannot read.
None of that cancels the present record. The grammar permits the record
to say: this much was carried, this much was reduced, this much remains,
and this is the time at which the reading stopped. It forbids the record
from claiming the authority of trials it did not perform.

Sequential recording illustrates the necessity. An image may be scanned
line by line, a spectrum sampled point by point, or a multidimensional
process serialized into a single ledger. The richness of the phenomenon
does not remove the order of commitment. A measuring act cannot commit
every distinction at once. It must place one recoverable distinction
after another. Later reconstruction may restore a surface, a field, or a
shape, but the ledger's authority comes from the ordered finite acts by
which the record was made.

This also explains why the horizon is not simply a practical limit. Even
a perfect instrument, if such a thing were imagined, would need a
grammar that says when a reading has been accepted. Without that
acceptance, there is no trial to repeat and no residue to compare. An
unending refinement may be a promise of knowledge, but it is not yet a
measurement. Measurement must stop somewhere in order to be read at all.

The stopping is not arbitrary when the grammar is doing its work. It is
bounded by the trial's rule of return. A clock supplies stages. A
stimulus supplies what counts as asking again. An expectation supplies
what counts as the same kind of answer. A residue supplies what remains
after the answer. The horizon gathers these into a finite statement:
under these conditions, by this much time, the ledger reads this.

The horizon also protects the distinction between unavailable and false.
A reading that cannot be obtained within the bound is not thereby
denied. It may be outside the current trial, outside the gauge, outside
the available time, or outside the permitted refinement. The grammar
forbids the cheap conversion of nonarrival into negation. This discipline
will matter when truth itself becomes the object of measurement. If the
grammar has learned not to treat an unavailable reading as false, it will
also have to face what happens when affirmation and denial have no
interior separator.

At the horizon, the apparatus is finite in the only way that matters for
measurement: it must commit a record before it has exhausted all
possible refinements. The grammar carries that finitude upward. It does
not apologize for it. It turns finitude into the condition under which
reading, repetition, and comparison can have force.
